\section{Data}
\label{sec:data}

\subsection{BBC News Corpus}
We use the BBC News dataset introduced by \citet{greene2006practical}, a widely-used benchmark for topic classification. The corpus contains 2,225 news articles collected from the BBC website between 2004 and 2005, spanning five editorial categories: \textit{business}, \textit{entertainment}, \textit{politics}, \textit{sport}, and \textit{tech}. The category distribution and exact split counts are shown in Table~\ref{tab:dataset}.

The articles are relatively short (median length approximately 300 words), topically diverse within each category, and written in a formal news register. Because BBC articles are carefully edited, the corpus has relatively few typos or non-standard forms, which simplifies preprocessing.

\subsection{Train / Validation / Test Split}

We apply a stratified split to preserve the class distribution in every partition. The dataset is divided into 80\% training, 10\% validation, and 10\% test, yielding 1,777 training, 221 validation, and 227 test articles. Stratification ensures that even the smaller categories (entertainment, tech) are adequately represented in every partition. The random seed is fixed at 42 for reproducibility.

\begin{table}[!htbp]
\centering
\small
\setlength{\tabcolsep}{3.8pt}
\begin{tabular}{lccccc}
\toprule
\textbf{Category} & \textbf{Articles} & \textbf{Share (\%)} & \textbf{Train} & \textbf{Val.} & \textbf{Test} \\
\midrule
Business      & 510 & 22.9 & 408 & 51 & 51 \\
Entertainment & 386 & 17.3 & 308 & 38 & 40 \\
Politics      & 417 & 18.7 & 333 & 41 & 43 \\
Sport         & 511 & 23.0 & 408 & 51 & 52 \\
Tech          & 401 & 18.0 & 320 & 40 & 41 \\
\midrule
\textbf{Total} & \textbf{2,225} & \textbf{100.0} & \textbf{1,777} & \textbf{221} & \textbf{227} \\
\bottomrule
\end{tabular}
\caption{BBC News dataset category distribution and exact stratified 80/10/10 split counts used for genre classification.}
\label{tab:dataset}
\end{table}

\subsection{Sentiment Labels}

No gold-standard sentiment labels exist for the BBC corpus. Instead, we generate pseudo-labels using VADER \cite{hutto2014vader}: articles with a compound score $\geq 0.05$ are labeled \textit{positive}, those $\leq -0.05$ are labeled \textit{negative}, and strictly neutral articles are discarded. These cutoffs follow VADER's standard guidance: scores close to zero are weak evidence of sentiment, and in factual news writing they are usually better treated as neutral. Discarding the borderline cases reduces label noise and produces a binary sentiment dataset drawn entirely from BBC text, keeping the vocabulary of the sentiment model consistent with the genre classifier's training domain.
